\documentclass[11pt]{article}
\usepackage{amsmath,amssymb,amsthm,geometry,hyperref}
\geometry{margin=1in}
\newtheorem{theorem}{Theorem}
\newtheorem{lemma}{Lemma}
\title{Identical stable tails do not determine an FI module\\A proof of Conjecture 00000005300}
\author{GitHub: gaochengzhecpu}
\date{}
\begin{document}
\maketitle

\section*{Statement and interpretation}
The source asks for two FI modules with all stable groups isomorphic but
nonisomorphic total modules, separated explicitly by their finite skeleton
before stabilization. We work over $\mathbb Z$. Here FI is the category of
finite sets and injective maps, and an FI module is a covariant functor
from FI to $\mathbb Z$-modules. An isomorphism of FI modules is a natural
isomorphism, not an unrelated collection of bijections.

The source does not specify a separate meaning for ``stable groups.'' We
prove the stronger concrete assertion that the restrictions to \emph{all
nonempty finite sets} are naturally isomorphic. Thus the example has
isomorphic eventual module values, compatible transition maps, and equal
stable direct limits. We do not interpret ``stable groups'' as an
unspecified additional invariant required to retain unstable information.

\begin{theorem}
There are finitely generated FI modules $T$ and $Z$ over $\mathbb Z$ such
that their restrictions to nonempty finite sets are naturally isomorphic,
but $T$ and $Z$ are not isomorphic as FI modules. Their only difference is
at the empty set, where $T(\varnothing)\cong\mathbb Z$ and
$Z(\varnothing)=0$.
\end{theorem}

\section*{Construction and proof}
Let $Z$ be the zero functor. Define
\[
 T(S)=\begin{cases}\mathbb Z,&S=\varnothing,\\0,&S\ne\varnothing.\end{cases}
\]
For an injection $f:S\hookrightarrow U$, use the identity map on
$\mathbb Z$ when both sets are empty and the unique zero map in every
other possible case. There is no injection from a nonempty set to the
empty set. Consequently identities are preserved. If the source and
target of a composable pair are empty, its intermediate set is empty
too, so all maps are identities. In every other case both the composite
map and the prescribed map of the composite are zero. This proves all
functor laws.

At every nonempty set both functors have value zero. The unique
isomorphisms between these zero modules commute with every injection
between nonempty sets, yielding a natural isomorphism on the entire
positive-size subcategory. In particular both standard stabilization
sequences are zero from size one onwards, so both direct limits are zero.

If a natural isomorphism $T\cong Z$ existed, its component at the empty
set would be an isomorphism $\mathbb Z\cong0$. This is impossible: an
injective homomorphism cannot identify $1$ and $0$. Thus the whole FI
modules are nonisomorphic. The element $1\in T(\varnothing)$ generates
$T$: its multiples span degree zero, and the images of those multiples
span the zero module in every positive degree. The zero module is
generated by the empty set of generators. The discrepancy is therefore
an explicit finite skeleton supported solely in size zero.

\section*{Formal verification and scope}
The Lean development constructs an actual category with objects $n\in
\mathbb N$ and morphisms $\operatorname{Fin}(n)\hookrightarrow
\operatorname{Fin}(m)$. This is the standard finite-set skeleton of FI.
Every finite set is isomorphic to one of its objects, so the construction
above on all finite sets has exactly this restriction. It also bundles
abelian groups with their $\mathbb Z$-module structures and takes
\emph{all} linear maps as morphisms; identities, composition and category
axioms are verified by Lean. The objects \texttt{supported} and
\texttt{zero} are Mathlib category-theoretic functors into this category.

For a definition without a dependent case split, Lean implements
\[
 T(n)=\mathbb Z^{E_n},\qquad
 E_n=\{\text{functions }\operatorname{Fin}(n)\longrightarrow\varnothing\}.
\]
The set $E_0$ is a singleton and $E_n$ is empty for $n\ge1$. Precomposition
with an injection gives $E_m\to E_n$, and precomposition again gives the
covariant linear map $T(n)\to T(m)$. This is exactly the functor above.
The theorem \texttt{degreeZeroEquiv} identifies $T(0)$ linearly with
$\mathbb Z$. Positive-size values are zero by
\begin{center}\texttt{supported\_subsingleton}.\end{center}
The definition \texttt{tailIso} is an actual natural
isomorphism on the positive-size subcategory, and \texttt{not\_iso}
excludes \emph{every} natural isomorphism of the full functors.
The theorem
\begin{center}\texttt{generated\_in\_degree\_zero}\end{center}
proves that every element is a multiple of the image of the displayed generator.

The main theorem \texttt{conjecture5300} proves an existential statement
over FI functors with both eventual linear equivalences and a natural
isomorphism of their positive restrictions, together with full
nonisomorphism. The direct-limit consequence and extension from the
skeleton to arbitrary finite sets are justified above on paper; Lean
does not construct categorical colimits or the equivalence with the
nonskeletal category. The kernel audit reports only \texttt{propext}
and \texttt{Quot.sound}; there is no placeholder proof, custom axiom, or
native decision procedure.

\section*{Source and reproduction}
The unchanged bilingual statement is included as \texttt{SOURCE.md},
from repository path
\begin{center}\texttt{conjectures/00000005300.md}.\end{center}
The source was checked at upstream commit
\begin{center}\small\texttt{b355e42906e6eaf52bc7e9d005e0ef02d8aa7400}.\end{center}
The project pins Lean 4.19.0 and Mathlib commit
\begin{center}\small\texttt{c44e0c8ee63ca166450922a373c7409c5d26b00b}.\end{center}
From \texttt{lean/}, run \texttt{lake build}, then
\begin{center}\small\texttt{lake env lean -DwarningAsError=true Main.lean}.\end{center}
The public Git manifest locks dependency revisions. The proof was
developed by a delegated AI agent with a self-review; the coordinating
agent reviews it separately. No external independent review is claimed.
\end{document}
